% New actual maps after S17. GPT-6 Astra Pro, /root/fixed_surface.
% 22 November 2026 UTC; effort not exposed.
\section[HD-F1-S18: norm-fourfold scalar action]{HD-F1-S18: scalar self-Ext and the exact holomorphic-form relation of the norm fourfold}
\label{sec:f1-norm-fourfold-scalar-action}

Retain the actual smooth connected norm fourfold $T$, its Abel map
$a:T\to P$, its reduced image $i:Y\hookrightarrow P$, and
$E=Ra_*\mathcal O_T=i_*\mathcal O_Y$ from S13--S17.
The following are computations on this specified object.
There is a canonical split injection
\[
 H^j(Y,\mathcal O_Y)=H^j(T,\mathcal O_T)
               \longrightarrow\operatorname{Ext}^j_P(E,E)
                       \quad(j\geq0).
 \tag{F1S18.1}
\]
The restrictions of ambient holomorphic forms to $T$ have ranks
\[
 \operatorname{rank}\bigl(H^0(P,\Omega_P^j)
                \xrightarrow{a^*}H^0(T,\Omega_T^j)\bigr)
                      =1,8,28,56,69\qquad(j=0,1,2,3,4).
 \tag{F1S18.2}
\]
In degree four the one-dimensional kernel is given explicitly below.
Its complex conjugate is an actual vanishing degree-four scalar
Hochschild action on $E$. The degree-two full trace is injective on
the 28-dimensional subspace in (F1S18.1) obtained by restricting
$H^2(P,\mathcal O_P)$. These statements do not compute the full
degree-two trace on $\operatorname{Ext}^2_P(E,E)$.

\paragraph{1. Canonical global maps giving the scalar retract.}
In the derived category of $Y$, form the actual derived self-intersection
\[
 \mathcal D=Li^*i_*\mathcal O_Y
          =\mathcal O_Y\otimes^{\mathbf L}_{\mathcal O_P}\mathcal O_Y.
\]
The first tensor factor gives its $\mathcal O_Y$-module structure.
The quotient unit $\eta:\mathcal O_P\to i_*\mathcal O_Y$ in the
second factor induces an $\mathcal O_Y$-linear map
\[
 u:\mathcal O_Y=\mathcal O_Y\otimes^{\mathbf L}_{\mathcal O_P}
                         \mathcal O_P\longrightarrow\mathcal D.
\]
Derived multiplication gives $\mu:\mathcal D\to\mathcal O_Y$.
The unit identity is the equality of actual derived morphisms
$\mu u=1_{\mathcal O_Y}$.
Derived adjunction now identifies
\[
 R\operatorname{Hom}_P(i_*\mathcal O_Y,i_*\mathcal O_Y)
       =R\operatorname{Hom}_Y(\mathcal D,\mathcal O_Y).
\]
Precomposition with $\mu$ gives a map from
$R\Gamma(Y,\mathcal O_Y)$ to this complex; precomposition with
$u$ is its retraction. Taking cohomology proves the split injection
(F1S18.1). The equality with $H^j(T,\mathcal O_T)$ follows from
the actual identity $Rb_*\mathcal O_T=\mathcal O_Y$, where
$a=ib$, proved in S16--S17. Thus no degeneration of a local-to-global
Ext spectral sequence, or global splitting of a chosen resolution,
is being presumed.

Under this inclusion, an ambient class
$\alpha\in H^j(P,\mathcal O_P)\subset HH^j(P)$ acts by its
ordinary restriction $i^*\alpha$. Indeed scalar multiplication on
the actual quotient algebra, followed by the unit, gives this
restriction, and the adjunction above gives its multiplication map
on $E$. The same maps on $T$ are $a^*\alpha$ followed by derived
pushforward: the lax monoidal multiplication on $Ra_*\mathcal O_T$
is the ordinary quotient-algebra multiplication under S17.

\paragraph{2. Positivity for forms of degree at most three.}
The actual class from S14 is
\[
             [Y]=\frac{\beta^4}{72}+w,\qquad
          \beta w=0,\qquad \int_Pw^2=18,
          \qquad\int_P\frac{\beta^8}{8!}=81.
 \tag{F1S18.3}
\]
For $1\leq j\leq3$ and a nonzero ambient holomorphic $j$-form
$\Omega$, the projection formula gives
\[
 i^{j^2}\int_Ta^*\Omega\,\overline{a^*\Omega}\,(a^*\beta)^{4-j}
       =\frac{i^{j^2}}{72}
                       \int_P\Omega\bar\Omega\,\beta^{8-j}>0.
 \tag{F1S18.4}
\]
The term involving $w$ vanishes because $4-j\geq1$ and
$\beta w=0$. The last inequality follows directly in a unitary
coframe for the positive form $\beta$: distinct exterior monomials
are orthogonal, and every diagonal coefficient is a positive
multiple of the ambient volume form. In particular $a^*\Omega$
cannot vanish. The ambient dimensions are $\binom8j$, proving
the first four entries of (F1S18.2), with the constant entry
following from connectedness of $T$.

\paragraph{3. The full degree-four Gram matrix and its actual kernel.}
Choose a $\beta$-unitary basis $a_1,\ldots,a_4,b_1,\ldots,b_4$
of the two holomorphic $K$-eigenspaces, and put
\[
 A=a_1\wedge\cdots\wedge a_4,\qquad
 B=b_1\wedge\cdots\wedge b_4,\qquad
 V=\int_PAB\bar A\bar B=81.
\]
The actual determinant plane in degree eight has complex basis
$A\bar B,B\bar A$. Since $w$ is real, write
\[
 w=\lambda A\bar B+\bar\lambda B\bar A.
\]
All reorderings of these degree-four factors have positive sign,
and hence
\[
         \int_Pw^2=2V|\lambda|^2=162|\lambda|^2=18,
                   \qquad |\lambda|=\frac13.
 \tag{F1S18.5}
\]
For the 70 exterior monomials in the chosen holomorphic basis,
consider $G_{IJ}=\int_Ta^*\Omega_I\,\overline{a^*\Omega_J}$.
The term $\beta^4/72$ gives the diagonal coefficient
$4!V/72=27$. The $w$ term is zero unless the two monomials are
$A,B$ in opposite positions. Thus the complete matrix is $27I$
on the other 68 monomials, and on $A,B$ its block is
\[
                   \begin{pmatrix}
                      27&81\bar\lambda\\
                      81\lambda&27
                    \end{pmatrix}.
 \tag{F1S18.6}
\]
It has eigenvalues zero and 54. The norm of the actual form
\[
                       \Omega_0=A-3\bar\lambda B
 \tag{F1S18.7}
\]
under this pairing is zero. A holomorphic top form on the smooth
projective fourfold $T$ has zero integral norm exactly when it
is zero. Therefore $a^*\Omega_0=0$, and (F1S18.6) proves that
this is the entire kernel. This proves the last entry of
(F1S18.2). The unitary change $b_1'=3\bar\lambda b_1$, with all
other basis elements fixed, gives $B'=3\bar\lambda B$ and sends
the coefficient $\lambda$ to $1/3$. It makes the same actual
kernel $\mathbb C(A-B')$; the phase change is this specified
basis map.

Complex conjugation and Hodge theory on $P,T$ give
\[
 \ker\bigl(H^4(P,\mathcal O_P)\to H^4(T,\mathcal O_T)\bigr)
       =\mathbb C\bigl(\bar A-3\lambda\bar B\bigr).
\]
Combining with the actual scalar maps of paragraph 1 proves
the relation on the original sheaf
\[
          \operatorname{ev}_E^4(\bar A-3\lambda\bar B)=0.
 \tag{F1S18.8}
\]
This is a computed characteristic-action relation, beyond a
relation obtained only by contracting the Chern character.

In fact these maps determine the entire scalar image as an algebra.
Put $W=H^1(P,\mathcal O_P)$ and
$\rho=\bar A-3\lambda\bar B\in\bigwedge^4W$.
The scalar action is the ring map induced by restriction and
multiplication in the quotient algebra. Since
$H^j(T,\mathcal O_T)=0$ for $j>4$, the split injection and all
the ranks just proved give the actual graded-algebra isomorphism
\[
 \operatorname{im}\bigl(H^*(P,\mathcal O_P)
                 \longrightarrow\operatorname{Ext}^*_P(E,E)\bigr)
       \simeq\frac{\bigwedge^*W}{(\rho,\,\bigwedge^5W)}.
 \tag{F1S18.8a}
\]
Its dimensions in degrees zero through four are
$(1,8,28,56,69)$, and its higher degrees vanish. In particular
every scalar quintic is an actual action relation, not only a
zero product with the Chern character.

There is also a geometric contraction map on the smooth locus of
$Y$. Since $\Omega_0|_Y=0$, contraction descends to
\[
             N_{Y/P}\longrightarrow\Omega_Y^3,
                   [v]\longmapsto(\iota_v\Omega_0)|_Y.
 \tag{F1S18.9}
\]
On the nonempty open set where $A|_Y\ne0$, both projections of
$T_yY$ to the two four-dimensional eigenspaces are isomorphisms;
the second one has determinant fixed by (F1S18.7). Write the
tangent plane as the graph of its actual invertible matrix $M$.
A normal vector represented by $(0,v)$ is sent by (F1S18.9) to
the nonzero scalar $-3\bar\lambda\det M$ times
$\iota_{M^{-1}v}A$. Thus (F1S18.9) is an isomorphism on that
specified open set. Nonemptiness follows from
$\int_Ta^*A\,\overline{a^*A}=27$. The formula does not identify
the normal sheaf across the zero divisor of $a^*A$ or across
the singularities of $Y$.

\paragraph{4. The degree-two scalar trace is actually injective.}
For $\alpha\in H^2(P,\mathcal O_P)$, the actual composition
of characteristic evaluation and the full Hochschild trace is
\[
         \sigma_E^2(\operatorname{ev}_E^2(\alpha))
                               =\alpha\operatorname{ch}(E).
 \tag{F1S18.10}
\]
These are precisely the diagonal-resolution characteristic map
and Atiyah--Hochschild trace of Buchweitz--Flenner,
Proposition 5.2.2, Proposition 6.2.1 and Corollary 6.3.2.
In this scalar summand the
contraction is ordinary cup product. The leading term on the
right is $\alpha[Y]\in H^6(P,\Omega_P^4)$. For nonzero $\alpha$,
the integral
\[
      \int_P\alpha\bar\alpha\,\beta^2[Y]
            =\frac1{72}\int_P\alpha\bar\alpha\,\beta^6>0
 \tag{F1S18.11}
\]
proves that this leading term is nonzero. Hence the full trace
is injective on these 28 actual independent scalar self-Ext
classes. Their independence also follows from (F1S18.1) and
the degree-two case of (F1S18.2).

The scalar retract is defined on all of $H^2(Y,\mathcal O_Y)$,
but its dimension and the full trace on the remaining scalar
classes have not been computed here. The other summands of
$\operatorname{Ext}^2_P(E,E)$ also remain to be computed. Thus
(F1S18.8) is an additional actual relation and (F1S18.10) is an
actual injective part of the trace, without a semiregularity
verdict for this eightfold sheaf.
